\subsection{Fase 1: comprensi\'on del negocio}
La tarea busca explicar el comportamiento del precio de vivienda \textit{medv} y su relaci\'on con variables estructurales, ambientales y socioecon\'omicas. El enfoque se limita a las fases 1 a 5 de CRISP-DM porque el producto esperado es un informe anal\'itico y no un sistema en producci\'on.

\subsection{Fase 2: comprensi\'on de los datos}
El conjunto Boston contiene 506 observaciones y 14 variables. No se detectaron valores faltantes, pero s\'i outliers en varias columnas, lo que confirma distribuciones asim\'etricas y colas largas. La Fig.~\ref{fig:outliers} resume este comportamiento.

\begin{figure}[htbp]
    \centering
    \includegraphics[width=0.9\textwidth]{boxplots_outliers.pdf}
    \caption{Detecci\'on visual de valores at\'ipicos en variables seleccionadas.}
    \label{fig:outliers}
\end{figure}

La matriz de correlaci\'on muestra relaciones fuertes entre variables urbanas y socioecon\'omicas, especialmente entre \textit{rad} y \textit{tax}, y una asociaci\'on negativa marcada entre \textit{lstat} y \textit{medv}. Esto sugiere que el valor de la vivienda est\'a vinculado a la condici\'on social del entorno.

\subsection{Fase 3: preparaci\'on de los datos}
Para el PCA se excluyeron \textit{chas} y \textit{medv} del bloque explicativo, porque la primera es una variable binaria contextual y la segunda es la variable de inter\'es principal. Luego se estandarizaron las 12 variables restantes para evitar que la escala dominara la soluci\'on.

\subsection{Fase 4: modelado PCA}
El scree plot confirma que el primer componente explica 51.06\% de la varianza y que cinco componentes acumulan 80\%. Esto permite reducir la dimensionalidad sin perder informaci\'on esencial.

La estructura de cargas indica que PC1 concentra principalmente variables asociadas a urbanizaci\'on y presi\'on urbana, como \textit{indus}, \textit{nox}, \textit{rad}, \textit{tax}, \textit{age} y \textit{lstat}, mientras que \textit{rm} y \textit{dis} se mueven en sentido contrario. En t\'erminos interpretativos, PC1 puede leerse como un gradiente de deterioro/urbanizaci\'on versus calidad residencial. Este resultado es consistente con la varianza explicada: PC1 aporta 51.06\% y cinco componentes concentran m\'as de 80\%.

\subsection{Fase 5: evaluaci\'on de las preguntas}
La pregunta sobre suburbios m\'as baratos se responde con el ranking de \textit{medv}: las observaciones 399, 406, 401, 400 y 415 est\'an entre las de menor valor, con precios entre 5.0 y 7.0. La Tabla~\ref{tab:respuestas} sintetiza las evidencias principales para cada pregunta.

\begin{table}[htbp]
\centering
\small
\begin{tabular}{p{2.0cm}p{4.9cm}p{5.0cm}}
\toprule
Pregunta & Evidencia principal & Conclusi\'on \\
\midrule
P1 & Outliers en al menos ocho variables, con mayor presencia en \textit{black}, \textit{zn} y \textit{crim}. & S\'i existen valores at\'ipicos. \\
P2 & Menores valores de \textit{medv} en las observaciones 399, 406, 401, 400 y 415. & Hay suburbios claramente m\'as baratos. \\
P3 & Correlaci\'on \textit{rm-medv} = 0.6954 y pendiente positiva significativa. & M\'as habitaciones implica mayor precio. \\
P4 & Media de \textit{medv}: 22.09 si \textit{chas}=0 y 28.44 si \textit{chas}=1; $p=0.003567$. & La cercan\'ia al r\'io eleva el valor. \\
P5 & Correlaci\'on \textit{lstat-medv} = -0.7377 y R$^2$ = 0.5441. & El estatus socioecon\'omico bajo reduce el valor. \\
P6 & Modelo con \textit{rad, tax, lstat, dis, medv, indus, nox} logra R$^2$ = 0.4337; con scores PCA y \textit{PC1--PC4}, R$^2$ = 0.6126. & La criminalidad es explicable en parte y el PCA ayuda a resumirla. \\
\bottomrule
\end{tabular}
\caption{S\'intesis de resultados para las preguntas de investigaci\'on.}
\label{tab:respuestas}
\end{table}

La relaci\'on entre tama\~no de la vivienda y precio es fuerte y positiva, como muestra la Fig.~\ref{fig:rmmedv}. Del mismo modo, la comparaci\'on por \textit{chas} evidencia diferencias estad\'isticamente significativas en \textit{medv}, por lo que la cercan\'ia al r\'io Charles no parece ser neutral para el mercado inmobiliario.

\begin{figure}[htbp]
    \centering
    \includegraphics[width=0.84\textwidth]{relacion_rm_medv.pdf}
    \caption{Relaci\'on entre el n\'umero medio de habitaciones y el valor de la vivienda.}
    \label{fig:rmmedv}
\end{figure}

\begin{figure}[htbp]
    \centering
    \includegraphics[width=0.82\textwidth]{boxplot_medv_chas.pdf}
    \caption{Valor de viviendas seg\'un cercan\'ia al r\'io Charles.}
    \label{fig:chas}
\end{figure}

En cuanto a \textit{lstat}, la relaci\'on negativa es una de las m\'as fuertes del conjunto y confirma que el nivel socioecon\'omico del entorno es un predictor relevante del precio. Finalmente, para \textit{crim} se observa que variables como \textit{rad}, \textit{dis}, \textit{lstat} y \textit{medv} aportan informaci\'on estad\'isticamente significativa, mientras que el modelo basado en componentes principales mejora la capacidad explicativa global.